\documentclass[a4paper,11pt]{article}
\usepackage[utf8]{inputenc}
\usepackage[T1]{fontenc}
\usepackage[english]{babel}
\usepackage[left=2.5cm,right=2.5cm,top=2cm,bottom=2cm]{geometry}
\usepackage{hyperref}
\usepackage{xcolor}
\usepackage{microtype}
\usepackage{lmodern}

\definecolor{primary}{RGB}{35, 55, 123}
\hypersetup{colorlinks=true, urlcolor=primary}

\begin{document}
\noindent \textbf{Lo\"ic Aron MBASSI EWOLO} \\
ENSEA -- Double Dipl\^ome Ing\'enieur \\
\href{mailto:loicaron.mbassiewolo@ensea.fr}{loicaron.mbassiewolo@ensea.fr} \quad (+33) 7 59 52 33 98 \\
\href{https://github.com/Nameless0l}{github.com/Nameless0l} \quad \href{https://mbassiloic.tech}{mbassiloic.tech}

\vspace{1em}\noindent Cergy, le \today
\vspace{1em}\noindent \textbf{Objet : Candidature Alternance -- Développeur IA, Vision par Ordinateur \& Computer Vision}
\vspace{1.5em}\noindent Madame, Monsieur,

\vspace{0.8em}\noindent J'ai d\'evelopp\'e un mod\`ele Computer Vision pour d\'etection d'anomalies cardiaques par analyse d'ECG photograph\'e dans le cadre du programme UROP 2024 (collaboration Google DeepMind), avec entra\^{\i}nement sur datasets h\'et\'erog\`enes et d\'eploiement de pipelines d'inf\'erence optimis\'es. Formation double dipl\^ome ENSEA/Polytechnique Yaound\'e, sp\'ecialit\'e syst\`emes embarqu\'es et math\'ematiques appliqu\'ees. Actuellement en stage INRIA Saclay (projet PRIMO) sur l'analyse statique d'applications mobiles et pipelines Python avanc\'es.

\vspace{0.8em}\noindent Affluences, sp\'ecialiste de la mesure d'affluence temps r\'eel (Louvre, Tour Eiffel, SNCF, Disney), repr\'esente un cadre id\'eal pour appliquer ces comp\'etences Computer Vision \`a la d\'etection et classification automatis\'ee de fr\'equentation. J'ai optimis\'e des mod\`eles TensorFlow Lite pour contraintes hardware h\'et\'erog\`enes (Arduino, Raspberry Pi) au Stage IoT Lab ENSPY, une exp\'erience directement transposable au d\'eploiement de syst\`emes de vision embarqu\'es sur vos capteurs.

\vspace{0.8em}\noindent Mon exp\'erience Urban Waste Management (1\textsuperscript{er} prix CITS National hackathon) a valid\'e ma maîtrise des architectures IoT complet, acquisition capteurs, ML pr\'edictif (scikit-learn, XGBoost) et \'evaluation mod\`eles \`a grande \'echelle. J'ai aussi d\'evelopp\'e un syst\`eme complet de reconnaissance faciale temps r\'eel (ClassSync) pour appel automatique, combinant OpenCV, Python et optimisation performance multiples cam\'eras en parall\`ele.

\vspace{0.8em}\noindent Je suis disponible en alternance d\`es septembre 2026 pour une dur\'ee de 12 mois, mobile en Île-de-France. Je serais ravi de discuter de comment mes comp\'etences Computer Vision, optimisation d'inf\'erence embarqu\'ee et pipeline ML avanc\'e peuvent contribuer aux innovations d'Affluences en mesure d'affluence autonome.

\vspace{1.5em}\noindent Veuillez agr\'eer, Madame, Monsieur, l'expression de mes salutations distingu\'ees.
\vspace{2em}\noindent \textbf{Lo\"ic Aron MBASSI EWOLO}
\end{document}
